\documentclass[10pt]{article}
\usepackage{amsmath}
\usepackage{amssymb}
\usepackage{amsthm}
\usepackage{braket}
\usepackage{hyperref}
\usepackage{setspace}

\doublespacing
\allowdisplaybreaks

\begin{document}
\section*{Discretization and Storage}
To store an arbitrary 1d state on qubits we do two things: we first restrict the domain of our wavefunction,
considering it on an interval $[-d, d]$ for $d > 0$, and we also discretize the
interval at steps $\Delta x > 0$. Now our state, rather than being of the form
$\ket{\psi(t)} = \int_{\mathbb{R}} \psi(x, t) \ket{x} dx$, is a normalized
Riemann sum $\ket{\hat{\psi}(t)} = \sum_{k = -\frac{d}{\Delta x}}^{\frac{d}{\Delta x}} a_k \ket{k \Delta x}$
($a_k \in \mathbb{C}$). We let $N = \frac{2d}{\Delta x}+1$ denote the number of grid positions.
This can be stored on $n = \lceil \log_2\left(N\right) \rceil = O\left(\log \left(\frac{d}{\Delta x} 
\right) \right)$ qubits with the computational basis state $\ket{k}$ representing
the position $-d + k \Delta x$, where $k$ is interpreted as an unsigned integer.
\section*{Free Particle}
In general, we are trying to solve $i \partial_t \psi = H \psi$, where 
$H = -\partial_x^2 + V(x)$. In the zero potential case, this simplifies to the PDE
$i \partial_t \psi(x, t) = -\partial_x^2 \psi(x, t)$. Applying the Fourier transform
$\mathcal{F}$ in x to both sides, we get $i \partial_t \mathcal{F}(\psi)(\xi, t) = \xi^2 \mathcal{F}(\psi)(\xi, t)$.
This gives the solution in Fourier space $\mathcal{F}(\psi)(\xi, t) = e^{-i \xi^2}\mathcal{F}(\psi)(\xi, 0)$.
Applying the inverse Fourier transform, we have $\psi(x, t) = \mathcal{F}^{-1}\left(e^{-i \xi^2} \mathcal{F}(\psi)(\xi, 0)\right)$.
For the discretization on qubits, we use the quantum Fourier transform, mapping
$\ket{k} \mapsto QFT e^{-i \xi_k^2 t} QFT^\dag \ket{k}$. Note that we apply
the inverse quantum Fourier transform first; this is because, due to unfortunate
choice of convention, $QFT^\dag: \ket{k} \mapsto \frac{1}{\sqrt{N}} \sum_{j=0}^{N-1} e^\frac{-2\pi ijk}{N} \ket{k}$
is the appropriate choice of approximation for $\mathcal{F}(\psi)(\xi, t) = \frac{1}{\sqrt{2\pi}}\int_{\mathbb{R}} e^{-i \xi x} \psi(x, t) dx$.
The choice of $\xi_k$ is the only decision here I can't really justify, but from \href{https://arxiv.org/pdf/1307.3318}{this paper},
we take $\Delta \xi = \frac{\pi}{d}$ and $\xi_0 = -\frac{\pi N}{2d}$, hence $\xi_k = -\frac{\pi N}{2d} + k \Delta \xi 
= \frac{\pi}{2d}(2k - N)$. 
\newline
We consider the quantum circuit implementation of $e^{-i \xi_k^2}$. Note that the
operator $A$ such that $e^A \ket{k} = e^k \ket{k}$ is given by
\begin{align*}
    A &= \sum_{j=0}^{n=1} \frac{I - Z_j}{2}2^j \text{\quad where $Z_j$ is the Pauli $Z$ on the $j$th qubit} \\
      &= \frac{1}{2} \sum_{j=0}^{n-1} 2^j(I - Z_j) \\
      &= \frac{1}{2}\left((2^n-1)I - \sum_{j=0}^{n-1} 2^j Z_j\right)
\end{align*}
Now that we have an idea of implementing $e^k$, consider $2k - N = (2^n - 1)I - \sum_{j=0}^{n-1} 2^j Z_j - NI
= -(I + \sum_{j=0}^{n-1}2^j Z_j)$, since $N = 2^n$. Hence,
\begin{align*}
    (2k-N)^2 &= I^2 + 2 \sum_{j=0}^{n-1} 2^j Z_j + \left( \sum_{j=0}^{n-1} 2^j Z_j \right) \left( \sum_{l=0}^{n-1} 2^l Z_l \right) \\
             &= I + \sum_{j=0}^{n-1} 2^{j+1} Z_j + \sum_{j=0}^{n-1} \sum_{l=0}^{n-1} 2^{j+l} Z_j Z_l \\
             &= I + \sum_{j=0}^{n-1} 2^{j+1} Z_j + \sum_{j=0}^{n-1} \sum_{l=j}^{n-1} 2^{j+l+1} Z_j Z_l \text{\quad since $Z_j Z_l = Z_l Z_j$}
\end{align*}
Finally we derive the gate decomposition for this operation. Letting $L = 2d$,
\begin{align*}
    e^{-i \xi_k^2 \Delta t} &= e^{-i \left(\frac{\pi}{L}(2k-N)\right)^2 \Delta t} \\
                            &= \exp\left(-i \frac{\pi^2}{L^2} \Delta t\left(I + \sum_{j=0}^{n-1} 2^{j+1} Z_j + \sum_{j=0}^{n-1} \sum_{l=j}^{n-1} 2^{j+l+1} Z_j Z_l\right)\right) \\
                            &\stackrel{(*)}{=} \left(\prod_{j=0}^{n-1} \exp\left( \frac{-i \pi^2 \Delta t}{L^2} 2^{j+1} Z_j \right)\right) \left( \prod_{j=0}^{n-1} \prod_{l=j+1}^{n-1} \exp\left(
                            \frac{-i \pi^2 \Delta t}{L^2} 2^{j+l+1} Z_j Z_l \right)\right) \\
                            &= \prod_{j=0}^{n-1} R_{Z_j}\left(2^{j+2} \frac{-i \pi^2}{L^2} \Delta t\right) \prod_{j=0}^{n-1} \prod_{l=j+1}^{n-1} R_{jl}\left(2^{j+l+2} \frac{-i \pi^2}{L^2} \Delta t\right) 
                            \text{\quad by definition}
\end{align*}
$(*)$ up to global phase \newline
$R_z$ and $R_{zz}$ gates are easily accessible on a quantum computer. This is the
circuit returned by kinetic() of 1d-solver.py.
\section*{Quantum Harmonic Oscillator}
As before, we are trying to solve $i \partial_t \psi = H \psi$, where
$H = \frac{p^2}{2m} + V$, where $p = -i \hbar \partial_x$ and $V$ is the
potential term. However, instead of $V \equiv 0$, we now consider the quantum
harmonic oscillator, with $V(x) = \frac{1}{2} m \omega^2 x^2$. For simplicity
of our units, we take $m = \frac{1}{2}$ (also do this in the previous section), $\hbar = 1$, and $\omega = 2$, thus hoping to solve
the PDE $i \partial_t \psi = -\partial_x^2 \psi + x^2 \psi$. After discretization 
of $\psi(x, 0)$ onto $n$ qubits, with $N = 2^n$ grid positions, we use the Trotter method, implementing $e^{-i p^2 \Delta t}
e^{-i V(x) \Delta t}$ for $\Delta t$ small and iterating to our desired time $t$.
Then $x_k^2 = (-d + k \Delta x)^2 = d^2 - 2d k \Delta x + k^2 \Delta x^2$, hence
$e^{-i x_k^2 \Delta t} \overset{(*)}{=} e^{-i(-2d k \Delta x \Delta t)} e^{-i k^2 \Delta x^2 \Delta t}$, where I
denote by $\overset{(*)}{=}$ equality up to global phase. Using the same decomposition
of $k = \sum_{j=0}^{n-1} \frac{I - Z_j}{2} 2^j$ as in the kinetic term previously,
\begin{align*}
e^{-i(-2dk \Delta x) \Delta t} &= \exp \left(2id \Delta x \Delta t \sum_{j=0}^{n-1}\frac{I-Z_j}{2} 2^j \right) \\
&= \exp \left(\sum_{j=0}^{n-1} i 2^j d \Delta x \Delta t(I-Z_j)\right) \\
&\overset{(*)}{=} \exp \left(\sum_{j=0}^{n-1} -i 2^j d \Delta x \Delta t Z_j \right)
\end{align*}
Considering now the $k^2$ term,
\begin{align*}
    k^2 &= \left(\sum_{j=0}^{n-1}\frac{I - Z_j}{2} 2^j \right)^2 \\
    &= \frac{1}{4} \left(\sum_{j=0}^{n-1} (I - Z_j)2^j \right) \left( \sum_{l=0}^{n-1} (I - Z_l)2^l \right) \\
    &= \frac{1}{4} \sum_{j=0}^{n-1} \sum_{l=0}^{n-1} (I-Z_j)(I-Z_l)2^{j+l} \\
    &\overset{(*)}{=} \frac{1}{4} \sum_{j=0}^{n-1} \sum_{l=0}^{n-1} (Z_j Z_l - Z_j - Z_l)2^{j+l} \\
    &= \frac{1}{4}\left( \sum_{j=0}^{n-1} \sum_{l=0}^{n-1} 2^{j+l} Z_j Z_l - \sum_{j=0}^{n-1} \sum_{l=0}^{n-1} 2^{j+l} Z_j 
    - \sum_{j=0}^{n-1} \sum_{l=0}^{n-1} 2^{j+l} Z_l \right) \\
    &\overset{(*)}{=} \frac{1}{2} \left( \sum_{j=0}^{n-1} \sum_{l=j+1}^{n-1} 2^{j+l} Z_j Z_l 
    - \sum_{j=0}^{n-1} \sum_{l=0}^{n-1} 2^{j+l} Z_j \right) \\
    &= \sum_{j=0}^{n-1} \sum_{l=j+1}^{n-1} 2^{j+l-1} Z_j Z_l - \sum_{j=0}^{n-1} 2^{j-1}Z_j\sum_{l=0}^{n-1} 2^l \\
    &= \sum_{j=0}^{n-1} \sum_{l=j+1}^{n-1} 2^{j+l-1} Z_j Z_l - \sum_{j=0}^{n-1} 2^{j-1}(2^n - 1)Z_j
\end{align*}
Thus we can derive the quantum gate decomposition for $e^{-i V(x) \Delta t}: \ket{k} \mapsto e^{-i x_k^2 \Delta t} \ket{k}$
as follows:
\begin{align*}
    e^{-i V(x_k)} &\overset{(*)}{=} e^{-i(-2dk \Delta x) \Delta t} e^{-i k^2 \Delta x^2 \Delta t} \\
    &= \exp\left(\sum_{j=0}^{n-1} -i 2^j d \Delta x \Delta t Z_j \right) 
    \exp\left( -i \Delta x^2 \Delta t \left( \sum_{j=0}^{n-1} \sum_{l=j+1}^{n-1} 2^{j+l-1} Z_j Z_l - \sum_{j=0}^{n-1} 2^{j-1}(2^n - 1)Z_j \right) \right) \\
    &= \exp\left( \sum_{j=0}^{n-1} -i \Delta x \Delta t 2^{j-1}(2d + \Delta x(1 - 2^n)) Z_j \right) \exp\left(-i \Delta x^2 \Delta t \sum_{j=0}^{n-1} \sum_{l=j+1}^{n-1} 2^{j+l-1} Z_j Z_l \right) \\
    &= \prod_{j=0}^{n-1}\left( R_{Z_j}\left(2^j \Delta x \Delta t(2d + \Delta x(1-N))\right) \prod_{l=j+1}^{n-1} R_{Z_{jl}} \left( 2^{j+l} \Delta x^2 \Delta t \right) \right)
\end{align*}
Note that all of the operators we are concerned with in implementing the potential
are simultaneously diagonalized, so it is valid to split exponentials of sums, as
well as noting that gate order does not matter. 
\section*{Analytical Solutions}
The analytical solutions used to verify the correctness of this algorithm are sourced
from \href{https://arxiv.org/pdf/quant-ph/0307046}{this paper} (qho) and Journal of Computational Physics 188 (2003) 157–17
(free particle).
\end{document}
